\documentclass[11pt]{article}
\usepackage[margin=1in]{geometry}
\usepackage{amsmath}
\usepackage{algorithm}
\usepackage{algpseudocode}

\title{Pseudocode for Data Cleaning Step}
\author{}
\date{}

\begin{document}

\maketitle

\begin{algorithm}
\caption{Crop Feature Engineering}
\begin{algorithmic}[1]
\State $D_{\text{noNA}} \gets$ remove all rows in $D_{\text{soil}}$ that contain missing values
\State $D_{\text{noOutlier}} \gets$ remove outliers in $D_{\text{noNA}}$ using Z-score filtering
\State $D_{\text{clean}} \gets \pi_{N, P, K, \text{pH}}(D_{\text{noOutlier}})$ \Comment{Project onto the soil-related features}
\State $(D_{\text{train}}, D_{\text{val}}, D_{\text{test}}) \gets \Call{randomSplit}{D_{\text{clean}}, 0.8, 0.1, 0.1}$
\end{algorithmic}
\end{algorithm}

\begin{table}[h]
\centering
\caption{Summary of crop data cleaning results}
\begin{tabular}{|l|r|}
\hline
\textbf{Metric} & \textbf{Value} \\
\hline
Original rows & 2200 \\
Rows removed for missing values & 0 \\
Rows removed as outliers ($|z| > 3.00$) & 171 \\
Final cleaned rows & 2029 \\
\hline
\end{tabular}
\end{table}

\end{document}
